\begin{center}
    \Large{\textbf{Аннотация}}
\end{center}

Работа посвящена количественному анализу стабилизации поверхности эмпирической функции потерь параметрических моделей при увеличении объема обучающей выборки. Целью является получение оценок достаточного размера выборки на основе локальной геометрии оптимизационной поверхности.

Сформулирован унифицированный критерий $p$-сходимости~$\Delta_p$ с функцией предпочтения точек в пространстве параметров, охватывающий одноточечную оценку и изотропное среднеквадратичное усреднение. Доказана верхняя оценка~$\mathcal O(k^{-2})$ для подпространственного критерия~$\Delta_2^{(D)}$ в подпространстве ведущих собственных векторов матрицы Гессе, в которой зависимость от размерности пространства параметров заменена на зависимость от размерности подпространства~$D$. Установлены замкнутое спектральное представление критерия и экстремальность подпространства главных направлений. Разработаны масштабируемые алгоритмы оценивания на основе произведений матрицы Гессе на вектор и моментов гауссовских квадратичных форм.

Полученные оценки проверены на полносвязной нейронной сети на MNIST и на трансформерной языковой модели \texttt{nanochat d6}. Наблюдаемая скорость убывания критериев согласуется с теоретическим показателем~$\alpha = 2$; подпространственный критерий согласуется с полнопространственным при~$D/N < 10^{-6}$ в локальном режиме квадратичной аппроксимации. Результаты применимы для обоснованного выбора достаточного объема обучающей выборки в нейронных сетях современных размеров.
